\chapter{The Schema Graph}
\label{ch:schema-graph}

The \texttt{SchemaGraph}\index{SchemaGraph} models a database schema as a graph where tables are nodes and foreign key relationships are edges.  It powers the \texttt{QueryBuilder}'s automatic join resolution (Chapter~\ref{ch:query-builder}).  Most applications use it indirectly through \texttt{db.newQueryBuilder()}, but understanding it helps when debugging join paths or working with schemas that lack declared foreign keys.

\section{Graph Model}
\label{sec:sg-model}

A \texttt{SchemaGraph} is an undirected graph:

\begin{itemize}
  \item \textbf{Nodes} are database table names.
  \item \textbf{Edges} are foreign key relationships, each connecting two tables via their key columns.
  \item Edges are traversable in both directions: knowing that \texttt{employee.department\_id} references \texttt{department.department\_id} lets the builder join from either table to the other.
  \item Both single-column and composite (multi-column) foreign keys are supported.
\end{itemize}

\section{Building a Schema Graph}
\label{sec:sg-building}

\subsection{From Database Metadata}

The simplest approach reads the schema directly from the live database connection using JDBC \texttt{DatabaseMetaData}:

\begin{lstlisting}[style=java]
SchemaGraph graph = SchemaGraph.fromDatabase(conn);
\end{lstlisting}

\noindent
This discovers all tables and their foreign key relationships.  It works well for databases with properly declared foreign keys.

\subsection{Cached per Schema (Thread-Safe)}

For applications with multiple schemas or concurrent requests, use the cached variant:

\begin{lstlisting}[style=java]
SchemaGraph graph = SchemaGraph.fromDatabase(conn, "mySchema");
\end{lstlisting}

\noindent
This uses a thread-safe cache (double-checked locking with \texttt{ConcurrentHashMap}) so the schema graph is built only once per schema name, regardless of how many threads request it simultaneously.

\subsection{Programmatic Declaration}

For databases without declared foreign keys (common in legacy systems), build the graph manually:

\begin{lstlisting}[style=java]
SchemaGraph graph = new SchemaGraph();
graph.addForeignKey("employee", "department_id",
                    "department", "department_id");
graph.addForeignKey("assignment", "employee_id",
                    "employee", "employee_id");
graph.addForeignKey("assignment", "project_id",
                    "project", "project_id");
\end{lstlisting}

\subsection{Composite Foreign Keys}

\begin{lstlisting}[style=java]
graph.addForeignKey("order_line",
    new String[]{"order_id", "product_id"},
    "order_product",
    new String[]{"order_id", "product_id"});
\end{lstlisting}

\subsection{Hybrid Approach}

Start from database metadata and add undeclared relationships:

\begin{lstlisting}[style=java]
SchemaGraph graph = SchemaGraph.fromDatabase(conn);
// Add a relationship for a legacy table without declared FKs
graph.addForeignKey("audit_log", "user_id",
                    "employee", "employee_id");
\end{lstlisting}

\begin{tipbox}
The hybrid approach is the best option for databases that have some declared foreign keys but also contain legacy tables without formal constraints.
\end{tipbox}

\section{File Caching}
\label{sec:sg-caching}

Building a schema graph from database metadata takes a few hundred milliseconds for large schemas.  For faster startup, save the graph to a file and load it later:

\begin{lstlisting}[style=java]
// Save (once, during setup or schema migration)
graph.saveToFile("schema-cache.txt");

// Load (sub-millisecond)
SchemaGraph graph = SchemaGraph.loadFromFile(
    "schema-cache.txt");
\end{lstlisting}

\noindent
The file format is a simple text format with \texttt{TABLE} and \texttt{FK} lines:

\begin{lstlisting}
TABLE employee
TABLE department
FK employee department_id department department_id
FK assignment employee_id employee employee_id
\end{lstlisting}

\section{In-Memory Cache Management}
\label{sec:sg-cache-mgmt}

The per-schema cache can be cleared when the database schema changes:

\begin{lstlisting}[style=java]
// Clear all cached schema graphs
SchemaGraph.clearSchemaCache();

// Clear a specific schema
SchemaGraph.clearSchemaCache("mySchema");
\end{lstlisting}

\section{Querying the Graph}
\label{sec:sg-querying}

\begin{lstlisting}[style=java]
// Check if a table exists in the graph
boolean exists = graph.hasTable("employee");

// Get all table names
Set<String> tables = graph.getTables();

// Get all edges (FK relationships) for a table
List<SchemaGraph.Edge> edges = graph.getEdges("employee");
for (SchemaGraph.Edge e : edges) {
    System.out.println(e.getFromTable() + "." +
        e.getFromColumn() + " -> " +
        e.getToTable() + "." + e.getToColumn());
}
\end{lstlisting}

\section{Finding Join Paths}
\label{sec:sg-paths}

The core algorithm finds the shortest join path connecting a set of tables:

\begin{lstlisting}[style=java]
Set<String> tables = Set.of("employee", "department",
    "building");
List<SchemaGraph.Edge> path =
    graph.findJoinPath(tables, "employee");
\end{lstlisting}

\noindent
The second parameter is the root table (the FROM table in the generated SQL).  The method returns an ordered list of edges representing the required JOINs.

The algorithm uses BFS (breadth-first search), which:
\begin{itemize}
  \item Always finds the shortest path (fewest joins).
  \item Runs in sub-millisecond time even for schemas with 300+ tables.
  \item Uses a greedy Steiner tree approximation for connecting multiple tables.
  \item Throws \texttt{SQLException} if no path exists between some tables.
\end{itemize}

\section{The Edge Inner Class}
\label{sec:sg-edge}

\texttt{SchemaGraph.Edge}\index{SchemaGraph.Edge} represents a foreign key relationship.  Key methods:

\begin{longtable}{P{4.8cm}P{5.7cm}}
\toprule
\textbf{Method} & \textbf{Description} \\
\midrule
\endhead
\texttt{getFromTable()} & Child table containing the FK column(s). \\
\texttt{getToTable()} & Parent/referenced table. \\
\texttt{getFromColumn()} & FK column name (first column if composite). \\
\texttt{getToColumn()} & Referenced column name (first if composite). \\
\texttt{getFromColumns()} & All FK column names (unmodifiable list). \\
\texttt{getToColumns()} & All referenced column names (unmodifiable list). \\
\texttt{isComposite()} & \texttt{true} if multi-column foreign key. \\
\texttt{buildOnCondition(from, to)} & Builds the ON clause string for a JOIN. \\
\bottomrule
\end{longtable}

\section{Integration with Connection}
\label{sec:sg-connection}

A \texttt{Connection} lazily builds and caches its schema graph:

\begin{lstlisting}[style=java]
// Get or build the schema graph for this connection
SchemaGraph graph = db.getSchemaGraph();

// Set the default schema
db.setSchema("myschema");
\end{lstlisting}

\noindent
When you call \texttt{db.newQueryBuilder()}, the builder automatically uses the connection's schema graph.  Most applications never interact with \texttt{SchemaGraph} directly.

\section{Integration with Command}
\label{sec:sg-command}

The \texttt{Command} class also provides \texttt{newQueryBuilder()}, which creates a \texttt{QueryBuilder} bound to that specific command.  The schema graph is obtained from the command's parent connection.

This is essential when multiple queries must be active simultaneously on the same connection.  Each \texttt{Command} wraps its own \texttt{PreparedStatement}, avoiding interference between concurrent result sets.  See Section~\ref{sec:qb-command} for usage examples.

\section{Example: Inspecting a Schema}
\label{sec:sg-example}

\begin{lstlisting}[style=groovy]
void inspectSchema(JSONObject injson, JSONObject outjson,
                   Connection db, ProcessServlet servlet) {
    SchemaGraph graph = db.getSchemaGraph()

    // List all tables
    Set<String> tables = graph.getTables()
    outjson.put("tableCount", tables.size())

    // For each table, list its relationships
    JSONArray tableInfo = new JSONArray()
    for (String table : tables) {
        JSONObject info = new JSONObject()
        info.put("name", table)
        List<SchemaGraph.Edge> edges = graph.getEdges(table)
        JSONArray rels = new JSONArray()
        for (SchemaGraph.Edge e : edges) {
            rels.put(e.getFromTable() + "." +
                e.getFromColumn() + " -> " +
                e.getToTable() + "." + e.getToColumn())
        }
        info.put("relationships", rels)
        tableInfo.put(info)
    }
    outjson.put("tables", tableInfo)
}
\end{lstlisting}
